% ---------------------------------------------------------------------------
% Experiment History -- CONDENSED, conference version (6-page IEEEtran).
% Replaces the whole \section{Experiment History} in main.tex (L534-) and the
% per-part drafts 1b, 5, 6, 7, 8, 9, 10 (cut 2026-09-15; recover the long
% versions from git, commit 4f6a1d5, if the thesis needs them).
%
% Budget: ~600 words of prose = ~1.1 IEEE columns, no table, no figure.
% It summarises; experiments that carry a claim are written out in full in
% their own sections, so no number here needs a table to be read.
%
% Carries the corrections README 3.5 lists as owed on the Overleaf parts 2-4:
%   - first experiment had 2 legitimate TX->RX pairs, not 1;
%   - two-agent numbers are job 99211's: BER ~0.24, per-agent power ~0.62,
%     det 2-8%, kurt ~-1.25 (NOT 0.33 / 0.45 / 3-10%); "double the single
%     agent" is gone -- the gain at matched total power is ~15%;
%   - the centralised-optimiser caveat is stated.
% Other numbers from: jobs 101306/101314/101328 (learned detector),
% 101622-101817 (raw-IQ RL), 101860/101866/102115 (characterisation),
% 101870/102305/102316/102319/102390 (realistic channel).
% ---------------------------------------------------------------------------

\section{Experiment History}
\label{sec:exphist}
We arrived at the current model by extending a simple link one step at a time.
This section summarises what each step established and why we moved on.
Unless stated otherwise, an attacker maximised $R = \mathrm{BER} - \beta D - \gamma P$, with $D$ a
detection penalty and $P$ the jammer power, trained either by PPO or by backpropagating directly
through a differentiable QPSK link built in Sionna.

\subsection{Statistical detectors}
A barrage jammer against two legitimate pairs on a lossless channel validated the chain: BER, power and
detection moved when, and only when, the jammer did.
A PPO jammer with a Gaussian policy then beat a power threshold ($\mathrm{BER} \approx 0.19$ at
$10\%$ detection) by transmitting just below it.
A kurtosis detector stopped it completely, since a Gaussian policy cannot produce the sub-Gaussian
statistics of QPSK, and neither a spline flow nor a Gaussian mixture trained by PPO did better.
The same flow trained by direct gradient evaded at once (kurtosis $-1.25$ against a threshold of
$-1.0$, BER up to $0.17$), but only with a white-box receiver and a genie view of the transmitted
symbol.

\textit{Two cooperating jammers.}
Trained jointly, two jammers reached $\mathrm{BER} \approx 0.24$ at $2$--$8\%$ detection and power
$0.62$ each.
The received constellation converged to four clean QPSK-like clusters with a quarter of the bits
wrong: the team hid by making the sum look legitimate.
At matched total power the gain over one jammer was about $15\%$, and with a single optimiser over
both agents this shows that a coordinated solution exists, not that decentralised agents find it.

\textit{Learned detection.}
The spectrogram CNN of Li \emph{et al.}~\cite{9707819} did not replicate on the flat carrier
($78.9\%$ accuracy), and caught barrage noise only at its $2.5\%$ false-alarm rate: there, a jammed
spectrogram is a clean one at lower SNR.
On an 802.11a-like OFDM link it replicated at once ($99.79\%$ accuracy, $0\%$ false alarms).

\textit{Policy gradient over raw IQ.}
MAPPO jammers emitting a complex value per subcarrier never learned against the frozen CNN, whose
output stayed at $0.999$ in every run.
Every broadband action the policy sampled was flagged, so the penalty was the same across the batch
and the normalised advantages carried no signal.
Reducing the action to two dimensions solved stealth but not damage ($\mathrm{BER}$ stuck at $0.35$),
and without the genie symbol cancellation is impossible for i.i.d.\ QPSK.
The obstacle was the action space, not the reward.

\textit{Detector characterisation.}
Sweeping attack parameters through the frozen CNN instead of training against it, interference at
equal power and equal damage ($\mathrm{BER} \approx 0.36$) was flagged $1.1\%$ of the time in-band and
$99.9\%$ once it also touched the empty guard and DC bins.
The replicated detector was a spectral-mask monitor.
Retraining on in-band jamming closed the gap at a price of $3.8\%$ false alarms, and adding a simple
energy detector cut the best stealthy BER from $0.42$ to $0.005$, because on a lossless channel any
added power is visible.

\textit{Realistic channel.}
With per-link frequency-selective fading and zero-forcing equalisation, a sparse in-band jammer imposed
a BER floor of $0.05$--$0.07$ that did not fall with SNR, and fading gave it cover again.
Two of our own numbers did not survive the right comparison.
A $70\%$ gain from channel-aware subcarrier selection, measured at equal configuration, fell to between
$-4.6\%$ and $+6.2\%$ at equal detectability.
And with the detector's own false-alarm rate as the stealth budget, no effective unflagged attacker
existed below $20$\,dB, and the best BER at $20$ and $30$\,dB was $0.011$ and $0.004$.

These steps fixed three choices for what follows: attacks are parameterised in low dimension rather
than raw IQ, results are compared at matched detectability under the false-alarm budget, and a learned
detector is judged against the optimal one rather than by its accuracy.
